\documentclass[12pt,a4paper]{article}
\usepackage[margin=1in]{geometry}
\usepackage{amsmath,amssymb}
\usepackage{booktabs}
\usepackage{array}
\usepackage{xcolor}
\usepackage{hyperref}
\usepackage{titlesec}
\usepackage{parskip}
\usepackage{graphicx}
\usepackage{float}
\usepackage{fancyhdr}
\usepackage{enumitem}
\usepackage{listings}
\usepackage{caption}

\definecolor{codeblue}{rgb}{0.1,0.2,0.6}
\definecolor{lightgray}{rgb}{0.95,0.95,0.95}
\definecolor{darkgray}{rgb}{0.3,0.3,0.3}
\definecolor{accentcolor}{rgb}{0.15,0.35,0.65}

\lstset{
    backgroundcolor=\color{lightgray},
    basicstyle=\ttfamily\small,
    keywordstyle=\color{codeblue}\bfseries,
    commentstyle=\color{darkgray}\itshape,
    breaklines=true,
    frame=single,
    rulecolor=\color{accentcolor},
    numbers=left,
    numberstyle=\tiny\color{darkgray},
    language=Python
}

\hypersetup{
    colorlinks=true,
    linkcolor=accentcolor,
    urlcolor=accentcolor
}

\pagestyle{fancy}
\fancyhf{}
\rhead{\small Deep Learning Assignment 1 -- Part A}
\lhead{\small Pravin Patil}
\cfoot{\thepage}
\renewcommand{\headrulewidth}{0.4pt}

\titleformat{\section}{\large\bfseries\color{accentcolor}}{}{0em}{\thesection\quad}[\titlerule]
\titleformat{\subsection}{\normalsize\bfseries\color{darkgray}}{}{0em}{\thesubsection\quad}

\begin{document}

%------------------------------------------------------------
%  TITLE PAGE
%------------------------------------------------------------
\begin{titlepage}
\centering
\vspace*{2cm}
{\Huge\bfseries\color{accentcolor} Deep Learning Assignment 1}\\[0.5em]
{\Large\bfseries Part A: Convolutional Neural Networks on CIFAR-100}\\[2em]
\rule{0.6\linewidth}{0.8pt}\\[1.5em]
{\large\textbf{Student:} Pravin Patil}\\[0.5em]
{\large\textbf{Framework:} TensorFlow 2.19.0 / Keras}\\[0.5em]
{\large\textbf{Dataset:} CIFAR-100 (32$\times$32 RGB, 100 classes)}\\[0.5em]
{\large\textbf{Date:} May 2026}\\[2em]
\rule{0.6\linewidth}{0.4pt}\\[1.5em]
\begin{abstract}
\noindent
This report documents a systematic investigation of Convolutional Neural Networks (CNNs) applied to the CIFAR-100 image classification benchmark. Three experimental configurations are studied: (A1) a CNN trained from scratch, (A2) robustness evaluation under additive Gaussian noise, (A3) noise-augmented retraining to improve robustness, and (A4/A5) transfer learning via a frozen VGG16 backbone. Quantitative results are presented alongside a rigorous analysis of numerical methods, hyperparameter choices, and the accuracy–robustness trade-off inherent in each approach.
\end{abstract}
\vfill
\end{titlepage}

%------------------------------------------------------------
%  1. INTRODUCTION
%------------------------------------------------------------
\section{Introduction}

Image classification at scale remains one of the central problems of computer vision. CIFAR-100~\cite{cifar} provides a challenging 100-class benchmark with 60{,}000 colour images of size $32\times32$ pixels ($50{,}000$ training, $10{,}000$ test). The low resolution forces a model to learn highly discriminative features from limited spatial information.

This assignment explores three model families:
\begin{enumerate}[noitemsep]
    \item A \textbf{custom CNN trained from scratch} with progressively increasing filter depth.
    \item \textbf{Noise-augmented retraining} to improve robustness against Gaussian perturbations.
    \item \textbf{Transfer learning} using VGG16 pre-trained on ImageNet as a frozen feature extractor.
\end{enumerate}

The unifying theme is the \emph{accuracy--robustness trade-off}: high accuracy on clean data does not imply resilience to distribution shift, and the three approaches occupy different positions along this spectrum.

%------------------------------------------------------------
%  2. DATASET AND PRE-PROCESSING
%------------------------------------------------------------
\section{Dataset and Pre-Processing}

\subsection{Data Statistics}

\begin{table}[H]
\centering
\caption{CIFAR-100 Dataset Summary}
\begin{tabular}{lrrrr}
\toprule
\textbf{Split} & \textbf{Samples} & \textbf{Image Size} & \textbf{Channels} & \textbf{Classes} \\
\midrule
Training & 50{,}000 & $32 \times 32$ & 3 (RGB) & 100 \\
Test     & 10{,}000 & $32 \times 32$ & 3 (RGB) & 100 \\
\bottomrule
\end{tabular}
\end{table}

\subsection{Normalisation}

Raw pixel values $p \in \{0, \ldots, 255\}$ are mapped to $[0,1]$:
\begin{equation}
    \hat{p} = \frac{p}{255}, \quad \hat{p} \in [0, 1].
\end{equation}
This affine rescaling ensures that the gradient magnitudes during back-propagation remain in a numerically stable range, preventing vanishing or exploding gradients in early layers.

\subsection{Label Encoding}

Integer class labels $y \in \{0, \ldots, 99\}$ are converted to one-hot vectors $\mathbf{y} \in \{0,1\}^{100}$ compatible with categorical cross-entropy loss. For label $c$:
\begin{equation}
    y_k = \begin{cases} 1 & k = c \\ 0 & k \neq c \end{cases}, \quad k \in \{0,\ldots,99\}.
\end{equation}

%------------------------------------------------------------
%  3. MODEL ARCHITECTURES
%------------------------------------------------------------
\section{Model Architectures}

\subsection{A1 -- Custom CNN from Scratch}

The baseline model follows a classic \emph{VGG-style} deepening strategy: the number of convolutional filters doubles at each block while spatial resolution is halved by max-pooling.

\begin{table}[H]
\centering
\caption{Custom CNN Layer Specification}
\begin{tabular}{llrr}
\toprule
\textbf{Layer} & \textbf{Type} & \textbf{Output Shape} & \textbf{Parameters} \\
\midrule
Input          & --           & $32\times32\times3$   & 0 \\
Conv2D-1       & Conv + ReLU  & $30\times30\times32$  & 896 \\
MaxPool-1      & MaxPool 2$\times$2 & $15\times15\times32$ & 0 \\
Conv2D-2       & Conv + ReLU  & $13\times13\times64$  & 18{,}496 \\
MaxPool-2      & MaxPool 2$\times$2 & $6\times6\times64$   & 0 \\
Conv2D-3       & Conv + ReLU  & $4\times4\times128$   & 73{,}856 \\
MaxPool-3      & MaxPool 2$\times$2 & $2\times2\times128$  & 0 \\
Flatten        & --           & 512                   & 0 \\
Dense-1        & FC + ReLU    & 128                   & 65{,}664 \\
Dropout(0.5)   & Regulariser  & 128                   & 0 \\
Dense-2        & FC + Softmax & 100                   & 12{,}900 \\
\midrule
\textbf{Total} & & & \textbf{171{,}812} \\
\bottomrule
\end{tabular}
\end{table}

\textbf{Design rationale.} Three convolutional blocks with filter counts $\{32, 64, 128\}$ extract low- to high-level features hierarchically. Each $3\times3$ kernel preserves local receptive field efficiency. Dropout ($p = 0.5$) on the penultimate layer acts as an ensemble method during training.

\subsection{A3 -- Noise-Augmented CNN}

The architecture is identical to A1. The key modification is in the training data: 20\% of training images are replaced by noise-corrupted copies, creating an augmented training set of 60{,}000 images.

\subsection{A4 -- Transfer Learning with VGG16}

VGG16 pre-trained on ImageNet is used as a frozen feature extractor. Because CIFAR-100 images are $32\times32$ -- far smaller than VGG16's typical $224\times224$ input -- each image is first up-sampled by a factor of $3$ to $96\times96$ via bilinear interpolation (UpSampling2D). A small classifier head ($256 \to 100$) is trained on top of the frozen convolutional features.

\begin{table}[H]
\centering
\caption{VGG16 Transfer Learning Head}
\begin{tabular}{llr}
\toprule
\textbf{Layer} & \textbf{Type} & \textbf{Trainable Parameters} \\
\midrule
UpSampling2D   & $3\times$ bilinear & 0 \\
VGG16 base     & Frozen conv blocks & 0 \\
Flatten        & --            & 0 \\
Dense(256)     & FC + ReLU     & 2{,}097{,}408 \\
Dropout(0.5)   & Regulariser   & 0 \\
Dense(100)     & FC + Softmax  & 25{,}700 \\
\midrule
\textbf{Total trainable} & & \textbf{2{,}123{,}108} \\
\bottomrule
\end{tabular}
\end{table}

%------------------------------------------------------------
%  4. NUMERICAL METHODS
%------------------------------------------------------------
\section{Numerical Methods}

\subsection{Loss Function}

All models minimise \textbf{categorical cross-entropy}:
\begin{equation}
    \mathcal{L}(\boldsymbol{\theta}) = -\frac{1}{N}\sum_{i=1}^{N}\sum_{k=0}^{99} y_{i,k} \log \hat{y}_{i,k},
\end{equation}
where $\hat{y}_{i,k} = \text{softmax}(\mathbf{z})_k$ is the predicted probability for class $k$ of sample $i$. The log ensures the loss surface is bounded below and has well-defined gradients even near extreme predictions.

\subsection{Optimisation Algorithm -- Adam}

All models use the \textbf{Adam} optimiser~\cite{adam}, which adapts per-parameter learning rates:
\begin{align}
    m_t &= \beta_1 m_{t-1} + (1-\beta_1)\,g_t, \\
    v_t &= \beta_2 v_{t-1} + (1-\beta_2)\,g_t^2, \\
    \hat{m}_t &= \frac{m_t}{1-\beta_1^t}, \quad \hat{v}_t = \frac{v_t}{1-\beta_2^t}, \\
    \theta_t &= \theta_{t-1} - \alpha \frac{\hat{m}_t}{\sqrt{\hat{v}_t}+\varepsilon},
\end{align}
where default hyperparameters $\alpha=10^{-3}$, $\beta_1=0.9$, $\beta_2=0.999$, $\varepsilon=10^{-7}$ are used. Adam's bias-correction terms ($1-\beta_1^t$, $1-\beta_2^t$) are critical in early epochs when the moment estimates are cold-started at zero.

\subsection{Activation Functions}

\textbf{ReLU} ($\max(0, x)$) is used in all hidden layers. It avoids the vanishing-gradient problem common with sigmoid/tanh activations in deep networks and is computationally trivial to differentiate. The \textbf{Softmax} output converts logits $\mathbf{z}$ to a valid probability distribution:
\begin{equation}
    \text{softmax}(\mathbf{z})_k = \frac{e^{z_k}}{\sum_{j=0}^{99} e^{z_j}}.
\end{equation}

\subsection{Convolution Operation}

For a 2D convolutional layer with kernel $\mathbf{W}$ of size $k \times k$ and $F$ filters applied to input feature map $\mathbf{X}$:
\begin{equation}
    \mathbf{Y}_{f,i,j} = \sum_{c}\sum_{m=0}^{k-1}\sum_{n=0}^{k-1} \mathbf{W}_{f,c,m,n}\,\mathbf{X}_{c,\,i+m,\,j+n} + b_f.
\end{equation}
All $3\times3$ kernels with unit stride (no padding) are used; each pooling layer halves spatial dimensions, preserving the most dominant activations.

\subsection{Gaussian Noise Model (A2)}

Additive Gaussian noise with variance $\sigma^2 = 0.05$ is applied:
\begin{equation}
    \tilde{x} = x + \epsilon, \quad \epsilon \sim \mathcal{N}(0, \sigma^2), \quad \sigma = \sqrt{0.05} \approx 0.2236,
\end{equation}
followed by clipping to $[0, 1]$. Given that $\hat{p} \in [0,1]$, the signal-to-noise ratio (SNR) is:
\begin{equation}
    \text{SNR} = \frac{\mathbb{E}[\hat{p}^2]}{\sigma^2} \approx \frac{0.25}{0.05} = 5,
\end{equation}
which represents a moderate perturbation -- enough to significantly degrade a model that has not been exposed to noise during training.

%------------------------------------------------------------
%  5. TRAINING CONFIGURATION
%------------------------------------------------------------
\section{Training Configuration and Hyperparameters}

\begin{table}[H]
\centering
\caption{Hyperparameter Summary Across All Models}
\begin{tabular}{lccc}
\toprule
\textbf{Hyperparameter} & \textbf{A1 (Scratch)} & \textbf{A3 (Aug.)} & \textbf{A4 (VGG16)} \\
\midrule
Epochs           & 80        & 80        & 80 \\
Batch size       & 128       & 128       & 64 \\
Optimiser        & Adam      & Adam      & Adam \\
Learning rate    & $10^{-3}$ & $10^{-3}$ & $10^{-3}$ \\
Dropout rate     & 0.5       & 0.5       & 0.5 \\
Training samples & 50{,}000  & 60{,}000  & 50{,}000 \\
Data augmentation & None     & Gaussian noise (20\%) & Rotation, shifts, flip \\
\bottomrule
\end{tabular}
\end{table}

\textbf{Batch size choice.} A batch size of 128 represents a compromise between gradient-estimate variance (smaller batches = noisier but regularising) and GPU throughput (larger batches = faster). VGG16 experiments use 64 due to the larger intermediate feature maps from the upsampled $96\times96$ inputs.

\textbf{Number of epochs.} Training for 80 epochs allows sufficient convergence of the A1 baseline without early stopping. Both A3 and A4 use the same epoch budget for a fair comparison.

%------------------------------------------------------------
%  6. RESULTS AND ANALYSIS
%------------------------------------------------------------
\section{Results and Analysis}

\subsection{Quantitative Performance}

\begin{table}[H]
\centering
\caption{Model Performance on CIFAR-100 -- Clean vs.\ Noisy Test Data}
\begin{tabular}{lcccc}
\toprule
\textbf{Model} & \textbf{Clean Acc.} & \textbf{Noisy Acc.} & \textbf{Abs. Drop} & \textbf{Rel. Drop} \\
\midrule
Scratch CNN (A1)           & 0.3825 & 0.0611 & 0.3214 & 84.07\% \\
Noise-Augmented CNN (A3)   & 0.3678 & 0.2440 & 0.1238 & 33.66\% \\
VGG16 Transfer (A4/A5)     & 0.4217 & 0.0266 & 0.3951 & 93.69\% \\
\bottomrule
\end{tabular}
\end{table}

\subsection{Observation -- Clean Data}

The VGG16 transfer model achieves the highest clean-data accuracy (42.17\%), benefiting from deep hierarchical feature representations learned on 1.2M ImageNet images. The scratch CNN reaches 38.25\% -- a respectable result for a lightweight model on a 100-class problem. The noise-augmented model (A3) trades a small 1.47 percentage point reduction in clean accuracy for drastically improved robustness.

\subsection{Observation -- Noisy Data}

Both A1 and A4 collapse under Gaussian noise. The scratch CNN drops to 6.11\% -- near random chance for 100 classes (expected random: 1\%). VGG16 fares even worse at 2.66\%, suggesting the ImageNet features are tuned to high-frequency texture patterns that are catastrophically disrupted by $\sigma = 0.224$ noise. The noise-augmented A3 model maintains 24.40\% accuracy, validating the strategy.

%------------------------------------------------------------
%  7. TRADE-OFF ANALYSIS
%------------------------------------------------------------
\section{Trade-Off Analysis}

\subsection{Accuracy vs.\ Robustness Trade-Off}

The core tension in this experiment is between \textbf{clean-data accuracy} and \textbf{noise robustness}. This is quantified by the \emph{robustness ratio}:
\begin{equation}
    \rho = \frac{\text{Acc}_{\text{noisy}}}{\text{Acc}_{\text{clean}}}.
\end{equation}

\begin{table}[H]
\centering
\caption{Robustness Ratio Comparison}
\begin{tabular}{lcc}
\toprule
\textbf{Model} & \textbf{Clean Acc.} & \textbf{Robustness Ratio $\rho$} \\
\midrule
Scratch CNN (A1)         & 0.3825 & 0.160 \\
Noise-Augmented CNN (A3) & 0.3678 & \textbf{0.663} \\
VGG16 Transfer (A4)      & 0.4217 & 0.063 \\
\bottomrule
\end{tabular}
\end{table}

A3 achieves a $\rho$ of 0.663, more than $4\times$ better than the scratch CNN and $10\times$ better than VGG16. This is obtained at a cost of only $\approx 1.5$ percentage points in clean accuracy.

\subsection{Model Complexity vs.\ Accuracy Trade-Off}

\begin{table}[H]
\centering
\caption{Model Complexity vs.\ Performance}
\begin{tabular}{lrrr}
\toprule
\textbf{Model} & \textbf{Parameters} & \textbf{Training Data} & \textbf{Clean Acc.} \\
\midrule
Scratch CNN (A1/A3) & 171{,}812   & 50{,}000 / 60{,}000 & 0.3825 / 0.3678 \\
VGG16 Transfer (A4) & 2{,}123{,}108 trainable & 50{,}000 & 0.4217 \\
\midrule
VGG16 total params  & $\approx$138M (frozen) & -- & -- \\
\bottomrule
\end{tabular}
\end{table}

VGG16 obtains +3.92 pp over A1 by leveraging 138M pretrained parameters, but its trainable head alone contains 12.4$\times$ more parameters than the entire scratch CNN. For embedded or edge deployments, A3 offers a better parameter-efficiency vs.\ robustness profile.

\subsection{Computational Cost Trade-Off}

\begin{table}[H]
\centering
\caption{Qualitative Computational Cost Comparison}
\begin{tabular}{lccc}
\toprule
\textbf{Model} & \textbf{Training Time} & \textbf{Inference Cost} & \textbf{Memory} \\
\midrule
Scratch CNN (A1)         & Low       & Low      & Low \\
Noise-Augmented CNN (A3) & Medium    & Low      & Low \\
VGG16 Transfer (A4)      & Medium    & High     & High \\
\bottomrule
\end{tabular}
\end{table}

The VGG16 approach requires upsampling every $32\times32$ image to $96\times96$ at inference time, tripling the spatial footprint and significantly increasing intermediate activation memory. The scratch models process the native resolution, making them 3--5$\times$ faster at inference.

\subsection{Data Augmentation Trade-Off (A3)}

Adding 20\% noisy copies increases training set size from 50{,}000 to 60{,}000 (a 20\% increase). The data augmentation introduces an inductive bias: the model learns that pixel-level Gaussian perturbations should not change the class label. This biases the decision boundary to be smooth in the $\ell_2$-ball around training points:
\begin{equation}
    f(\mathbf{x} + \boldsymbol{\epsilon}) \approx f(\mathbf{x}), \quad \|\boldsymbol{\epsilon}\|_2 \leq r,
\end{equation}
which is consistent with Lipschitz continuity of the learned function. The slight clean-accuracy drop (0.3825 $\to$ 0.3678) reflects the model allocating representational capacity to the noise manifold at the cost of fitting fine-grained clean-image features.

%------------------------------------------------------------
%  8. DISCUSSION
%------------------------------------------------------------
\section{Discussion}

\textbf{Why does VGG16 fail under noise?} ImageNet features learned by VGG16 encode high-frequency texture and colour statistics. Gaussian noise with $\sigma = 0.224$ is spectrally broad and heavily corrupts these high-frequency components, making the frozen feature representations unreliable for downstream classification. Furthermore, the upsampling step (UpSampling2D) creates block artefacts that interact poorly with noise, compounding the distribution shift.

\textbf{Why does A3 succeed?} Noise-augmented training is a form of \emph{data augmentation regularisation}. It implicitly trains the model to learn class-relevant low-frequency features that persist under additive noise, rather than over-fitting to texture patterns. This is related to adversarial training: exposing the model to perturbed inputs during training improves robustness to the same perturbation at test time.

\textbf{Limitations.} The improved robustness of A3 is noise-type-specific. A model trained against Gaussian noise may remain fragile to adversarial examples, blur, JPEG compression, or geometric distortions. A comprehensive robustness strategy would use diverse augmentation (random crops, colour jitter, CutMix, AugMix) alongside noise injection.

%------------------------------------------------------------
%  9. CONCLUSION
%------------------------------------------------------------
\section{Conclusion}

This report compared three CNN-based approaches to CIFAR-100 classification with a focus on robustness to additive Gaussian noise:

\begin{enumerate}[noitemsep]
    \item \textbf{Scratch CNN (A1)} achieves 38.25\% accuracy on clean data but degrades to 6.11\% under noise -- an 84\% relative drop.
    \item \textbf{Noise-Augmented CNN (A3)} reduces clean accuracy slightly to 36.78\% but maintains 24.40\% under noise, yielding a robustness ratio 4$\times$ better than A1.
    \item \textbf{VGG16 Transfer Learning (A4)} achieves the best clean-data accuracy (42.17\%) but is the most fragile under noise (2.66\%), performing near random-chance with a 93.7\% relative drop.
\end{enumerate}

The central finding is that \textbf{accuracy on clean data and robustness to noise are not automatically coupled}. Transfer learning from ImageNet -- despite its accuracy advantage -- introduces a strong prior over high-frequency features that makes it particularly sensitive to pixel-level perturbations. Noise-augmented training provides the best accuracy--robustness trade-off among the three strategies, at minimal additional computational cost.

Future work should explore Batch Normalisation in the scratch architecture, fine-tuning of the VGG16 backbone with noise-augmented data, and the use of certified defences such as randomised smoothing.

%------------------------------------------------------------
%  REFERENCES
%------------------------------------------------------------
\begin{thebibliography}{9}

\bibitem{cifar}
A.~Krizhevsky, \textit{Learning Multiple Layers of Features from Tiny Images}, Technical Report, University of Toronto, 2009.

\bibitem{adam}
D.~P.~Kingma and J.~Ba, \textit{Adam: A Method for Stochastic Optimization}, arXiv:1412.6980, 2014.

\bibitem{vgg}
K.~Simonyan and A.~Zisserman, \textit{Very Deep Convolutional Networks for Large-Scale Image Recognition}, ICLR 2015.

\bibitem{dropout}
N.~Srivastava et al., \textit{Dropout: A Simple Way to Prevent Neural Networks from Overfitting}, JMLR, 15(1):1929--1958, 2014.

\bibitem{relu}
V.~Nair and G.~E.~Hinton, \textit{Rectified Linear Units Improve Restricted Boltzmann Machines}, ICML 2010.

\end{thebibliography}

\end{document}
